\documentclass[11pt]{article}

% ---------- Packages ----------
\usepackage[margin=1in]{geometry}
\usepackage{setspace}        % for line spacing
\usepackage{amsmath,amssymb,amsfonts}
\usepackage{graphicx}
\usepackage{subcaption}
\usepackage{booktabs}
\usepackage{enumitem}
\usepackage{hyperref}
\usepackage[nameinlink]{cleveref}
\usepackage[numbers,sort&compress]{natbib}
\usepackage{url}

% ---------- Metadata ----------
\title{Title Goes Here:\\
  A VAR-based Computational Analysis of Influenza and Weather Dynamics}
\author{Your Name \and Collaborator Name}
\date{\today}

% ---------- Custom Commands ----------
\newcommand{\R}{\mathbb{R}}
\newcommand{\E}{\mathbb{E}}
\newcommand{\Var}{\mathrm{Var}}
\newcommand{\Cov}{\mathrm{Cov}}

\begin{document}
\onehalfspacing   % 1.5 spacing is nice for a 10+ page essay

\maketitle

\begin{abstract}
% TODO: 1--2 paragraphs.
% Briefly state:
% - Biological problem (e.g., influenza dynamics, weather, seasonality).
% - Computational formulation (VAR, Granger, IRFs, bootstrap).
% - Main dataset(s) used.
% - Key qualitative findings (e.g., which variables Granger-cause others, key IRF patterns).
\end{abstract}

\tableofcontents
\newpage

% ============================================================
\section{Introduction}
% ============================================================

\subsection{Biological Motivation and Problem Statement}
% TODO:
% - Introduce the influenza / epidemiological context.
% - Explain why understanding temporal dynamics and causal relationships matters
%   (e.g., forecasting, policy, intervention planning).
% - Define the scientific questions in words.

\subsection{From Biological Questions to Computational Problems}
% TODO:
% - Translate the biological questions into:
%   - Forecasting problems,
%   - Causal inference / Granger causality questions,
%   - Shock/response questions via IRFs.
% - Specify high-level tasks, e.g.:
%   - Predict log-transformed influenza counts from past values and weather.
%   - Test whether weather variables Granger-cause influenza incidence.
%   - Quantify the effect of a one-time shock in weather on future influenza.

\subsection{Overview of Approach}
% TODO:
% - Give a brief roadmap of the computational pipeline:
%   1. Data preprocessing and feature engineering.
%   2. Fitting a VAR($p$) model.
%   3. Computing impulse response functions (IRFs).
%   4. Performing Granger causality tests.
%   5. Bootstrapping IRFs and Granger statistics for uncertainty.
% - Mention software stack (Go backend for VAR/GC/IRF, R Shiny front-end, etc.).

% ============================================================
\section{Technical Background}
% ============================================================

\subsection{Time Series Analysis in Epidemiology}
% TODO:
% - Explain basic time series concepts: univariate vs multivariate, stationarity,
%   lag structures, seasonality, etc.
% - Discuss why multivariate modeling is needed (e.g., multiple weather variables + influenza).

\subsection{Vector Autoregression (VAR)}
% High-level explanation; detailed math goes to Appendix.
% Use intuitive language here.
% TODO:
% - Describe the idea: each variable depends on lagged values of all variables.
% - Introduce notation informally:
%   - $\mathbf{y}_t \in \R^K$ is the vector of variables at time $t$.
%   - VAR($p$) includes lags $1,\dots,p$.
% - Explain what ``reduced form'' means (all variables on left, error terms on right).

\subsection{Impulse Response Functions (IRFs)}
% TODO:
% - Intuition: IRFs trace out the effect of a one-time shock in one variable
%   on all variables over time.
% - Explain why IRFs are useful for interpretation in epidemiology/public health.

\subsection{Granger Causality}
% TODO:
% - Define Granger causality conceptually (prediction-based causality).
% - Give simple example: if past weather improves prediction of influena
%   beyond past influenza alone, then weather Granger-causes influenza.

\subsection{Bootstrap Methods for Time Series}
% TODO:
% - Explain why asymptotic standard errors might be inadequate (short samples, model misspecification).
% - Introduce residual bootstrap idea at a high level
%   (resample residuals, simulate new series, re-fit VAR, re-compute statistic).

% ============================================================
\section{Data and Preprocessing}
% ============================================================

\subsection{Datasets}
% TODO:
% - Describe influenza dataset (resolution, time span, variables, source).
% - Describe weather dataset (temperature, humidity, etc.).
% - Mention any external covariates (seasonality, sine/cosine terms, etc.).

\subsection{Preprocessing and Feature Engineering}
% TODO:
% - Resampling / aligning time series (e.g., weekly aggregation).
% - Handling missing data (interpolation, exclusion, imputation).
% - Transformations (log, differencing, scaling) applied to influenza counts and covariates.
% - Construction of seasonality regressors (e.g., sine/cosine).

\subsection{Final Variable Set}
% TODO:
% - Enumerate final variables $\mathbf{y}_t$ used in VAR.
% - Provide a small table listing variable name, meaning, units.

% Example table:
% \begin{table}[h]
%   \centering
%   \begin{tabular}{ll}
%     \toprule
%     Variable & Description \\
%     \midrule
%     LOG\_INF\_A & Log-transformed influenza A counts \\
%     TEMP       & Average weekly temperature (C) \\
%     HUMID      & Average weekly humidity (\%) \\
%     \dots      & \dots \\
%     \bottomrule
%   \end{tabular}
%   \caption{Variables included in the VAR model.}
%   \label{tab:variables}
% \end{table}

% ============================================================
\section{Computational Formulation}
% ============================================================

\subsection{Forecasting Problem}
% TODO:
% - Formulate forecasting as: given history $\{\mathbf{y}_{t-1},\dots,\mathbf{y}_{t-p}\}$,
%   predict $\mathbf{y}_t$.
% - Define performance metrics (e.g., RMSE, MAE, coverage).

\subsection{Causal and Structural Questions}
% TODO:
% - Formalize null hypotheses for Granger causality tests.
% - Explain how IRFs correspond to ``structural'' questions
%   (e.g., effect of weather shock on influenza trajectory).

\subsection{Model Selection and Assumptions}
% TODO:
% - State assumptions (e.g., linearity, covariance-stationary errors).
% - Describe lag-order selection (AIC/BIC/other).
% - Mention any stationarity checks / transformations.

% ============================================================
\section{Algorithms and Implementation}
% ============================================================

% Note: This section is high-level descriptions, not code dumps.

\subsection{VAR Estimation via OLS}
% TODO:
% - Explain stacking of lagged observations into a regression matrix.
% - Describe solving $(X^\top X)^{-1} X^\top Y$ or using SVD/pseudoinverse.
% - Mention handling of deterministic terms: constant, trend, seasonality.

\subsection{Impulse Response Computation}
% TODO:
% - Describe computing moving-average representation:
%   - Recursively compute $\Psi_h$ from $\{A_1,\dots,A_p\}$.
%   - Apply a structural shock vector (e.g., from Cholesky of $\Sigma_u$).
% - Explain choice of identification (e.g., Cholesky ordering) qualitatively.

\subsection{Granger Causality Testing}
% TODO:
% - Explain restricted vs unrestricted regressions for an effect variable.
% - Describe F-test construction:
%   - RSS difference,
%   - degrees of freedom,
%   - resulting $F$ statistic and $p$-value.
% - State significance level (e.g., $\alpha = 0.05$).

\subsection{Residual Bootstrapping for IRFs}
% TODO:
% - Explain pipeline:
%   1. Fit VAR once, compute residuals.
%   2. Resample residual vectors with replacement.
%   3. Simulate bootstrap series using fitted VAR and resampled residuals.
%   4. Re-fit VAR, re-compute IRFs.
%   5. Aggregate distribution across replications to form confidence bands.
% - Mention number of replications and rationale.

\subsection{Residual Bootstrapping for Granger Causality}
% TODO:
% - Similar pipeline:
%   1. For each bootstrap replication, simulate series.
%   2. Re-fit VAR and recompute Granger F-statistics.
%   3. Use empirical distribution to obtain bootstrap $p$-values.

\subsection{Software Architecture}
% TODO:
% - Describe your Go implementation as a modular library (TimeSeries, ReducedFormVAR, etc.).
% - Explain how R Shiny or other tools call the Go binary and visualize forecasts/IRFs/GC.
% - Include any parallelization or performance considerations (e.g., worker pools for bootstraps).

% ============================================================
\section{Results}
% ============================================================

\subsection{Model Fit and Forecasting Performance}
% TODO:
% - Present forecast error metrics on held-out data (tables/figures).
% - Show plots of observed vs forecasted influenza trajectories.

% Example figure skeleton:
% \begin{figure}[h]
%   \centering
%   \includegraphics[width=0.7\textwidth]{figs/forecast_plot.pdf}
%   \caption{Observed and forecasted influenza counts for Country X.}
%   \label{fig:forecast}
% \end{figure}

\subsection{Impulse Response Analysis}
% TODO:
% - Show IRF plots for key shocks (e.g., shock to weather variable, or influenza itself).
% - Interpret direction, magnitude, and duration of responses.

\subsection{Granger Causality Results}
% TODO:
% - Present Granger causality matrix (heatmap or table).
% - Highlight significant relationships (e.g., weather $\to$ influenza).

\subsection{Bootstrap Uncertainty Estimates}
% TODO:
% - Show IRF plots with bootstrap confidence bands.
% - Compare asymptotic vs bootstrap-based conclusions if relevant.
% - Show bootstrap-based Granger $p$-values and discuss robustness.

% ============================================================
\section{Discussion}
% ============================================================

% TODO:
% - Summarize main findings in biological terms (e.g., weather influences influenza with lag X).
% - Relate to prior work and literature.
% - Discuss strengths (multivariate, interpretable, uncertainty-quantified) and limitations
%   (linearity assumptions, data quality, nonstationarity, etc.).
% - Suggest extensions (e.g., SVAR, BVAR, non-linear VAR, other causal discovery methods).

% ============================================================
\section{Conclusion}
% ============================================================

% TODO:
% - Briefly restate the problem and main contributions.
% - Emphasize the link from first principles → computational model → real data insights.
% - Mention potential future directions (e.g., deployment, broader datasets, policy use).

% ============================================================
\section*{Acknowledgments}
% Optional: acknowledge mentors, collaborators, datasets, etc.

% ============================================================
\section*{References}
% ============================================================

% If you use BibTeX, delete this environment and just:
%   \bibliographystyle{plainnat}
%   \bibliography{refs}
% For now, a manual placeholder:

\begin{thebibliography}{99}

\bibitem{Hamilton1994}
J.~D. Hamilton.
\newblock \emph{Time Series Analysis}.
\newblock Princeton University Press, 1994.

% TODO: Add more references (epidemiology papers, Granger causality, bootstrap, etc.).

\end{thebibliography}

\newpage

% ============================================================
\appendix
% ============================================================

% ------------------------------------------------------------
\section{Supplementary Methods: Mathematical Details}
\label{app:math}
% ------------------------------------------------------------

In this appendix we have the mathematical details of the methods used:
vector autoregression (VAR), impulse response functions (IRFs), Granger
causality tests, and residual bootstrap procedures.

\subsection{Vector Autoregression (VAR)}
\label{app:var}
\subsubsection{Model specification}

Let $\{\mathbf{y}_t\}_{t=1}^T$ be a $K$-dimensional time series, where
$\mathbf{y}_t \in \mathbb{R}^K$ are all the variables (e.g., influenza
counts and weather covariates) at time $t$.
A VAR($p$) model with intercept can be written as
\begin{equation}
  \mathbf{y}_t
  = \mathbf{c}
  + \sum_{i=1}^p A_i \mathbf{y}_{t-i}
  + \mathbf{u}_t,
  \qquad t = 1,\dots,T,
  \label{eq:var-original}
\end{equation}
where
\begin{itemize}[leftmargin=1.5em]
  \item $\mathbf{c} \in \mathbb{R}^K$ is a vector of intercepts (can be modified to add deterministic constants)
  \item $A_i \in \mathbb{R}^{K \times K}$ are autoregressive coefficient matrices,
  \item $\mathbf{u}_t \in \mathbb{R}^K$ is the innovation (error) term.
\end{itemize}
We assume
\begin{equation}
  \mathbb{E}[\mathbf{u}_t] = \mathbf{0}, \qquad
  \mathbb{E}[\mathbf{u}_t \mathbf{u}_s^\top] =
    \begin{cases}
      \Sigma_u, & t = s, \\
      0, & t \neq s,
    \end{cases}
\end{equation}
so that $\{\mathbf{u}_t\}$ is simply random noise.

\subsubsection{Deterministic regressors}

We can write \cref{eq:var-original} with deterministic
regressors. Let $\mathbf{d}_t \in \mathbb{R}^M$ be a vector of deterministic
terms (e.g., constant, linear trend, seasonal dummies, sine/cosine pairs),
and let $B \in \mathbb{R}^{K \times M}$ be its coefficient matrix. Then
\begin{equation}
  \mathbf{y}_t
  = B \mathbf{d}_t
  + \sum_{i=1}^p A_i \mathbf{y}_{t-i}
  + \mathbf{u}_t.
  \label{eq:var-det}
\end{equation}
The pure-intercept case corresponds to $\mathbf{d}_t \equiv \mathbf{1}$ and
$M = 1$.

\subsubsection{Stacked regression form}

For estimation, we rewrite the VAR as a multivariate linear regression.

Define the \emph{effective sample size} $T^\star := T - p$ since the first few lags cannot be used. Now we can 
put all the depdenent variables from $t = p+1$ to $T$ into a matrix
\begin{equation}
  Y
  = \begin{bmatrix}
      \mathbf{y}_{p+1}^\top \\
      \mathbf{y}_{p+2}^\top \\
      \vdots \\
      \mathbf{y}_{T}^\top
    \end{bmatrix}
  \in \mathbb{R}^{T^\star \times K}.
\end{equation}
Next, define the regressor matrix
\begin{equation}
  X
  = \begin{bmatrix}
      \mathbf{x}_{p+1}^\top \\
      \mathbf{x}_{p+2}^\top \\
      \vdots \\
      \mathbf{x}_{T}^\top
    \end{bmatrix}
  \in \mathbb{R}^{T^\star \times (Kp + M)},
\end{equation}
where each row is
\begin{equation}
  \mathbf{x}_t^\top
  = \big[
      \mathbf{y}_{t-1}^\top,\,
      \mathbf{y}_{t-2}^\top,\,
      \dots,\,
      \mathbf{y}_{t-p}^\top,\,
      \mathbf{d}_t^\top
    \big]
  \in \mathbb{R}^{1 \times (Kp + M)}.
\end{equation}
We also stack the coefficient matrices into a single matrix
\begin{equation}
  B
  = \begin{bmatrix}
      A_1^\top \\
      A_2^\top \\
      \vdots \\
      A_p^\top \\
      B^\top
    \end{bmatrix}
  \in \mathbb{R}^{(Kp + M) \times K}.
\end{equation}
The innovations are stacked as
\begin{equation}
  U
  = \begin{bmatrix}
      \mathbf{u}_{p+1}^\top \\
      \mathbf{u}_{p+2}^\top \\
      \vdots \\
      \mathbf{u}_{T}^\top
    \end{bmatrix}
  \in \mathbb{R}^{T^\star \times K}.
\end{equation}
Then \cref{eq:var-det} for $t = p+1,\dots,T$ becomes the matrix equation
\begin{equation}
  Y = X B + U.
  \label{eq:var-stacked}
\end{equation}

\subsection{Estimation via Ordinary Least Squares}
The ordinary least squares (OLS) estimator of $B$ minimizes
$\|Y - X B\|_F^2$, the sum of squared residuals over all equations:
\begin{equation}
  \widehat{B}
  = \arg\min_{B} \|Y - X B\|_F^2.
\end{equation}
Assuming $X$ has full column rank, the solution is
\begin{equation}
  \widehat{B}
  = (X^\top X)^{-1} X^\top Y.
  \label{eq:ols-B}
\end{equation}
The fitted values and residuals are
\begin{equation}
  \widehat{Y} = X \widehat{B},
  \qquad
  \widehat{U} = Y - \widehat{Y}.
\end{equation}
Let $\widehat{\mathbf{u}}_t$ denote the row of $\widehat{U}$ corresponding
to time $t$.

An estimator of the innovation covariance matrix $\Sigma_u$ is
\begin{equation}
  \widehat{\Sigma}_u
  = \frac{1}{T^\star - (Kp + M)}
    \widehat{U}^\top \widehat{U}
  \in \mathbb{R}^{K \times K},
  \label{eq:sigma_u_hat}
\end{equation}
where the denominator uses the residual degrees of freedom.  Here $M$
denotes the number of deterministic regressors included in the VAR
(e.g., constant, trend, seasonal dummies, or other non-stochastic terms).


\subsubsection*{Rank-deficient regressors and the Moore--Penrose pseudoinverse}

If $X$ does \emph{not} have full column rank, then the matrix
$(X^\top X)$ is singular and the closed-form expression
\eqref{eq:ols-B} is not defined.  In this case the OLS estimator is
solved using the Moore--Penrose pseudoinverse $X^{+}$:
\begin{equation}
  \widehat{B}
  = X^{+} Y,
\end{equation}
which gives us the minimum-norm solution to the least-squares problem
$\min_{B} \|Y - X B\|_F^2$ even if $X$ is rank-deficient.  The
pseudoinverse is computed via singular value
decomposition (SVD).

\subsection{Moving-Average Representation and IRFs}
\label{app:irf}
\subsubsection{Companion form}

Let $Kp$-dimensional stacked state vector
\begin{equation}
  \mathbf{Y}_t
  = \begin{bmatrix}
      \mathbf{y}_t^\top &
      \mathbf{y}_{t-1}^\top &
      \dots &
      \mathbf{y}_{t-p+1}^\top
    \end{bmatrix}^\top
  \in \mathbb{R}^{Kp}.
\end{equation}
The VAR($p$) can be written in first-order form as
\begin{equation}
  \mathbf{Y}_t
  = \mathbf{C}
  + F \mathbf{Y}_{t-1}
  + \mathbf{U}_t,
  \label{eq:companion}
\end{equation}
where the \emph{companion matrix} $F \in \mathbb{R}^{Kp \times Kp}$ is
\begin{equation}
  F =
  \begin{bmatrix}
    A_1 & A_2 & \dots & A_{p-1} & A_p \\
    I_K & 0   & \dots & 0       & 0   \\
    0   & I_K & \dots & 0       & 0   \\
    \vdots & \vdots & \ddots & \vdots & \vdots \\
    0   & 0   & \dots & I_K     & 0
  \end{bmatrix},
\end{equation}
$\mathbf{C} \in \mathbb{R}^{Kp}$ contains the intercepts (and deterministic
terms), and
\begin{equation}
  \mathbf{U}_t =
  \begin{bmatrix}
    \mathbf{u}_t^\top &
    \mathbf{0}^\top &
    \dots &
    \mathbf{0}^\top
  \end{bmatrix}^\top
  \in \mathbb{R}^{Kp}.
\end{equation}

If the VAR is stable, then $(I_{Kp} - F)$ is invertible, and the mean is
\begin{equation}
  \mathbb{E}[\mathbf{Y}_t]
  = (I_{Kp} - F)^{-1} \mathbf{C}
  =: \boldsymbol{\mu}_Y.
\end{equation}
$\boldsymbol{\mu}_Y$.

\subsubsection{Infinite-order moving-average representation}

Under stability, we can rewrite the VAR in \emph{vector moving-average}
(VMA($\infty$)) form:
\begin{equation}
  \mathbf{y}_t
  = \boldsymbol{\mu}
  + \sum_{h=0}^{\infty} \Psi_h \mathbf{u}_{t-h},
  \label{eq:vma}
\end{equation}
where $\boldsymbol{\mu} \in \mathbb{R}^K$ is the unconditional mean of
$\mathbf{y}_t$ and $\Psi_h \in \mathbb{R}^{K \times K}$ are the VMA
coefficient matrices.
We set
\begin{equation}
  \Psi_0 = I_K,
\end{equation}
and find $\Psi_h$ recursively from the VAR coefficients:
\begin{equation}
  \Psi_h
  = A_1 \Psi_{h-1}
    + A_2 \Psi_{h-2}
    + \dots
    + A_p \Psi_{h-p},
  \qquad h \ge 1,
  \label{eq:psi-recursion}
\end{equation}
where $\Psi_h = 0$ for $h < 0$.

\subsubsection{Structural shocks and identification}

The reduced-form innovations $\mathbf{u}_t$ are usually correlated:
\begin{equation}
  \mathbb{E}[\mathbf{u}_t \mathbf{u}_t^\top] = \Sigma_u,
  \qquad \Sigma_u \text{ symmetric positive definite.}
\end{equation}
To interpret shocks as orthogonal or ``structural'' shocks, we factor
$\Sigma_u$ as
\begin{equation}
  \Sigma_u = P P^\top,
\end{equation}
where $P$ is a nonsingular $K \times K$ matrix.
We can do this using the Cholesky factorization with $P$ lower triangular,
which encodes some sort of ordering of the variables.
We can then write
\begin{equation}
  \mathbf{u}_t
  = P \boldsymbol{\varepsilon}_t,
  \qquad
  \boldsymbol{\varepsilon}_t \sim \mathcal{N}(\mathbf{0}, I_K),
\end{equation}
so that $\boldsymbol{\varepsilon}_t$ has uncorrelated unit-variance.
Substituting into \cref{eq:vma} gives
\begin{equation}
  \mathbf{y}_t
  = \boldsymbol{\mu}
  + \sum_{h=0}^{\infty} \Psi_h P \boldsymbol{\varepsilon}_{t-h}
  = \boldsymbol{\mu}
  + \sum_{h=0}^{\infty} \Theta_h \boldsymbol{\varepsilon}_{t-h},
\end{equation}
where
\begin{equation}
  \Theta_h := \Psi_h P \in \mathbb{R}^{K \times K}.
\end{equation}

\subsubsection{Impulse response functions}

Consider a one-unit shock to variable $j$ at time $t$, i.e.,
$\boldsymbol{\varepsilon}_t = \mathbf{e}_j$ (the $j$-th basis
vector) and $\boldsymbol{\varepsilon}_s = \mathbf{0}$ for $s \neq t$.
The effect on $\mathbf{y}_{t+h}$ for horizon $h \ge 0$ is
\begin{equation}
  \Delta \mathbf{y}_{t+h}
  = \Theta_h \mathbf{e}_j
  = \text{(the $j$-th column of $\Theta_h$)}.
\end{equation}
The \emph{impulse response function} from structural shock $j$ to variable $k$
at horizon $h$ is thus
\begin{equation}
  \mathrm{IRF}_{k \leftarrow j}(h)
  = \big[ \Theta_h \big]_{k,j}.
  \label{eq:irf-entry}
\end{equation}
In practice, we compute $\widehat{\Psi}_h$ from the estimated VAR
coefficients via \cref{eq:psi-recursion}, form
$\widehat{\Theta}_h = \widehat{\Psi}_h \widehat{P}$, and then find
$\widehat{\mathrm{IRF}}_{k \leftarrow j}(h)$ for $h=0,\dots,H$.

\subsection{Granger Causality Test Statistic}
\label{app:granger}
\subsubsection{Definition}

Let $y_{k,t}$ and $y_{j,t}$ be two components of the vector
$\mathbf{y}_t$.
Informally, $y_{j,t}$ \emph{Granger-causes} $y_{k,t}$ if past values of
$y_{j,t}$ contain information that helps predict $y_{k,t}$ beyond the
information contained in past values of $y_{k,t}$ and other variables.

More precisely, in the context of a VAR($p$), we compare two models for
the $k$-th equation:

\begin{itemize}[leftmargin=1.5em]
  \item \textbf{Unrestricted model (U):} includes lags of $y_{j,t}$:
  \begin{equation}
    y_{k,t}
    = \beta_{0}
    + \sum_{\ell=1}^K \sum_{i=1}^p \beta_{\ell,i}^{(U)}\, y_{\ell,t-i}
    + \boldsymbol{\gamma}^\top \mathbf{d}_t
    + u_{k,t}^{(U)}.
    \label{eq:gc-unrestricted}
  \end{equation}

  \item \textbf{Restricted model (R):} excludes all $p$ lags of $y_{j,t}$,
        i.e., imposes the linear restrictions
        $\beta_{j,1}^{(U)} = \dots = \beta_{j,p}^{(U)} = 0$:
  \begin{equation}
    y_{k,t}
    = \beta_{0}
    + \sum_{\ell=1}^K \sum_{i=1}^p \beta_{\ell,i}^{(R)}\, y_{\ell,t-i}
      \quad \text{with } \beta_{j,i}^{(R)} = 0 \text{ for all } i
    + \boldsymbol{\gamma}^\top \mathbf{d}_t
    + u_{k,t}^{(R)}.
    \label{eq:gc-restricted}
  \end{equation}
\end{itemize}

The \emph{null hypothesis} of no Granger causality from $y_{j,t}$ to
$y_{k,t}$ is
\begin{equation}
  H_0:
  \beta_{j,1}^{(U)} = \beta_{j,2}^{(U)} = \dots = \beta_{j,p}^{(U)} = 0,
\end{equation}
which says that lags of $y_{j,t}$ do not enter the $k$-th equation.

\subsubsection{F-statistic for linear restrictions}

Let $T^\star = T - p$ denote the effective sample size for the regression.
For the $k$-th equation, denote the OLS residuals from the unrestricted and
restricted regressions by $\widehat{u}_{k,t}^{(U)}$ and
$\widehat{u}_{k,t}^{(R)}$, respectively, and define the residual sums of
squares
\begin{equation}
  \mathrm{RSS}_U = \sum_{t=p+1}^T \big(\widehat{u}_{k,t}^{(U)}\big)^2,
  \qquad
  \mathrm{RSS}_R = \sum_{t=p+1}^T \big(\widehat{u}_{k,t}^{(R)}\big)^2.
\end{equation}
Let $q$ be the number of restrictions being tested; here $q = p$ because
we restrict $p$ lag coefficients of $y_{j,t}$ in the $k$-th equation.
Let $m_U$ be the number of parameters in the unrestricted equation
(including intercept and deterministic terms).
The usual $F$-statistic is
\begin{equation}
  F
  = \frac{
      (\mathrm{RSS}_R - \mathrm{RSS}_U) / q
    }{
      \mathrm{RSS}_U / (T^\star - m_U)
    }.
  \label{eq:F-stat}
\end{equation}
Under $H_0$ and standard regularity conditions, $F$ has an $F$-distribution
with $(q,\, T^\star - m_U)$ degrees of freedom:
\begin{equation}
  F \sim F_{q,\, T^\star - m_U}.
\end{equation}
We can then obtain a $p$-value by comparing the observed $F$ to this distribution,
and reject the null of no Granger causality if the $p$-value is less than
a chosen significance level (e.g., $\alpha = 0.05$).

In matrix form, multiple Granger tests can be performed simultaneously
using the estimated coefficient matrices $\widehat{A}_i$ and the estimated
innovation covariance $\widehat{\Sigma}_u$, but in our implementation we
apply \cref{eq:F-stat} equation-by-equation for the sake of interpretability.  

\subsection{Residual Bootstrap for IRFs}
\label{app:bootstrap-irf}


Analytic IRF confidence intervals are derived under assumptions such as
correct VAR specification, Gaussian and homoskedastic errors, and large-sample
asymptotics. In our application, these assumptions are not fully satisfied.
To obtain inference that is robust to such deviations, we use a residual
bootstrap, which approximates the sampling variability of the IRFs directly
from the observed data without relying on these parametric assumptions.


\subsubsection{Algorithm}

Let $\{\mathbf{y}_t\}_{t=1}^T$ be the observed series, and let
$\widehat{A}_1,\dots,\widehat{A}_p$, $\widehat{B}$, and
$\widehat{\Sigma}_u$ be the parameters estimated by OLS as described in
\cref{app:var}. Let $\widehat{\mathbf{u}}_t$ be the corresponding
residuals for $t = p+1,\dots,T$.

Now condsider $B$ bootstrap replications. For replication $b = 1,\dots,B$:

\begin{enumerate}[leftmargin=1.5em,label=\textbf{(B\arabic*)}]
  \item \textbf{Resample residuals.}
        Construct a bootstrap residual sequence
        $\{\mathbf{u}_t^{*(b)}\}_{t=p+1}^T$ by drawing with replacement from
        the set of centered residuals
        $\{\widehat{\mathbf{u}}_t - \bar{\mathbf{u}} : t=p+1,\dots,T\}$,
        where
        \begin{equation}
          \bar{\mathbf{u}}
          = \frac{1}{T^\star} \sum_{t=p+1}^T \widehat{\mathbf{u}}_t.
        \end{equation}
        Centering enforces zero mean in the bootstrap innovations. Because the
        distribution of the residuals is taken from the data rather than assumed,
        this step avoids parametric assumptions such as normality.

  \item \textbf{Initialize bootstrap series.}
        Choose starting values for the bootstrap series
        $\mathbf{y}_1^{*(b)},\dots,\mathbf{y}_p^{*(b)}$.
        Weset
        \begin{equation}
          \mathbf{y}_t^{*(b)} = \mathbf{y}_t,
          \qquad t = 1,\dots,p,
        \end{equation}
        Using the observed history ensures that each bootstrap sample shares the
        same starting conditions as the original dataset.

  \item \textbf{Simulate bootstrap series.}
        For $t = p+1,\dots,T$, generate
        \begin{equation}
          \mathbf{y}_t^{*(b)}
          = \widehat{B}\,\mathbf{d}_t
          + \sum_{i=1}^p \widehat{A}_i \mathbf{y}_{t-i}^{*(b)}
          + \mathbf{u}_t^{*(b)}.
          \label{eq:bootstrap-sim}
        \end{equation}
        This produces a bootstrap sample $\{\mathbf{y}_t^{*(b)}\}_{t=1}^T$.

  \item \textbf{Re-fit VAR to bootstrap sample.}
        Estimate a VAR($p$) on $\{\mathbf{y}_t^{*(b)}\}$ in the exact same
        way as for the original data, obtaining
        $\widehat{A}_i^{*(b)}$, $\widehat{B}^{*(b)}$ and
        $\widehat{\Sigma}_u^{*(b)}$.

  \item \textbf{Compute bootstrap IRFs.}
        From $\{\widehat{A}_i^{*(b)}\}$ and $\widehat{\Sigma}_u^{*(b)}$,
        compute the bootstrap VMA coefficients
        $\widehat{\Psi}_h^{*(b)}$ for $h=0,\dots,H$ using
        \cref{eq:psi-recursion}, obtain a bootstrap Cholesky factor
        $\widehat{P}^{*(b)}$ such that
        $\widehat{\Sigma}_u^{*(b)} = \widehat{P}^{*(b)}
          \widehat{P}^{*(b)\top}$,
        form $\widehat{\Theta}_h^{*(b)} =
          \widehat{\Psi}_h^{*(b)} \widehat{P}^{*(b)}$,
        and extract bootstrap IRFs
        \begin{equation}
          \widehat{\mathrm{IRF}}_{k \leftarrow j}^{*(b)}(h)
          = \big[ \widehat{\Theta}_h^{*(b)} \big]_{k,j},
          \qquad h = 0,\dots,H.
        \end{equation}
\end{enumerate}

After $B$ replications, we get distributions for each IRF
entry and horizon:
\begin{equation}
  \Big\{
    \widehat{\mathrm{IRF}}_{k \leftarrow j}^{*(b)}(h)
  \Big\}_{b=1}^B.
\end{equation}
A $(1-\alpha)$ confidence interval at horizon $h$ can be gotten
using quantiles, e.g. the percentile method:
\begin{equation}
  \Big[
    q_{\alpha/2},
    \;q_{1-\alpha/2}
  \Big],
\end{equation}
where $q_\tau$ is the $\tau$-quantile of the bootstrap sample. Because
these intervals are derived from the empirical distribution rather than
asymptotic approximations, they remain valid even when standard VAR assumptions
(such as normality or homoskedasticity) are violated.  

\subsection{Residual Bootstrap for Granger Causality}  
\label{app:bootstrap-gc}

We also use a residual bootstrap to obtain empirical distributions and
$p$-values for the Granger causality $F$-statistics in
\cref{eq:F-stat}.
The goal is to simulate many bootstrap series using the fitted VAR model,
re-fit the model to each bootstrap series, recompute the $F$-statistics, and
compare the observed statistic to this bootstrap distribution.

Let $F_{\text{obs}}$ be the observed $F$-statistic for a given
Granger test (e.g., from variable $j$ to $k$).
Using the same bootstrap series
$\{\mathbf{y}_t^{*(b)}\}_{t=1}^T$ generated in
\cref{app:bootstrap-irf}, we do:

\begin{enumerate}[leftmargin=1.5em,label=\textbf{(G\arabic*)}]
  \item \textbf{Re-fit VAR to bootstrap series.}
        For replication $b$, estimate the VAR($p$) on
        $\{\mathbf{y}_t^{*(b)}\}$, obtaining
        $\{\widehat{A}_i^{*(b)}\}$ and $\widehat{B}^{*(b)}$.

  \item \textbf{Compute bootstrap $F$-statistic.}
        For the Granger relationship of interest (e.g., $j \to k$),
        re-estimate the unrestricted and restricted versions of the
        $k$-th equation on the bootstrap series and compute
        $F^{*(b)}$ using the same formula as in \cref{eq:F-stat}.
\end{enumerate}

This yields bootstrap statistics
\begin{equation}
  \{ F^{*(b)} \}_{b=1}^B.
\end{equation}
A \emph{bootstrap $p$-value} for the observed $F_{\text{obs}}$ is then
\begin{equation}
  \widehat{p}_{\text{boot}}
  = \frac{
      1 + \#\{b : F^{*(b)} \ge F_{\text{obs}}\}
    }{
      B + 1
    }.
  \label{eq:bootstrap-pvalue}
\end{equation}
If $\widehat{p}_{\text{boot}} < \alpha$ (e.g., $\alpha = 0.05$), we regard
the Granger relationship as statistically significant at level $\alpha$,
based on the bootstrap approximation to the finite-sample distribution.

The bootstrap framework in \cref{app:bootstrap-irf,app:bootstrap-gc} thus
provides a unified way to quantify uncertainty in both impulse responses
and Granger causality tests, using the same core VAR model and residual
structure.

% ------------------------------------------------------------
\section{Additional Figures and Tables}
\label{app:figs}
% ------------------------------------------------------------

% TODO:
% - Place any extra plots, robustness checks, or tables that do not fit in the main text.

% Example:
% \begin{figure}[h]
%   \centering
%   \includegraphics[width=0.8\textwidth]{figs/extra_irf.pdf}
%   \caption{Additional IRF results for alternative lag order.}
%   \label{fig:extra_irf}
% \end{figure}

% ------------------------------------------------------------
\section{Author Contributions}
\label{app:contrib}
% ------------------------------------------------------------

% TODO: Fill this section at the end.
% Be explicit about who did what for the essay, code, and presentation.

% Suggested structure:
% - Person A: 
%   - Designed and implemented VAR estimation and forecasting code in Go.
%   - Implemented residual bootstrapping for IRFs and Granger causality.
%   - Wrote Sections X--Y and Appendix \ref{app:math}.
%
% - Person B:
%   - Collected and cleaned influenza and weather datasets.
%   - Developed R Shiny app for visualization and user interface.
%   - Wrote Sections Z--W and prepared figures.
%
% - Both:
%   - Co-designed the overall project.
%   - Prepared and delivered the project presentation.

\end{document}
